\documentclass[11pt]{article}
\usepackage[margin=1in]{geometry}
\usepackage{amsmath,amssymb,amsthm,mathtools,hyperref}
\DeclareMathOperator{\rank}{rank}
\DeclareMathOperator{\diag}{diag}
\title{Quantitative return of original mixed-action singular values}
\author{}
\date{20 September 2026}
\begin{document}
\maketitle
This is a receiving calculation from the proved constant-depth original
observation theorem DO1--23 and its differentiated-conductor extension
DC1--35. Both complete proofs accompany this note. All metrics are the
original attained metric and its actual observation quotient; the stated
finite constants are computed from those matrices. No unidentified
isometry, change of period or bound independent of degree is used.

\section{Original data and the proved recovery constant}
Retain the actual simple-quartet period locus of DO, $k\equiv1\pmod4$,
$k\ge77$, $q=(k+1)^2$, $r=8k-16$, the original polynomial quotient
$E_k$, its full attained metric $G_N$, and original multiplication $A=M_S$.
For the observation onto its image $\Lambda:E_k\to B_k$, put
\[
 Q_N=(\Lambda G_N^{-1}\Lambda^*)^{-1},\quad
 I:\mathbb C^r\overset\sim\longrightarrow\ker\Lambda,
 \quad H_K=I^*G_NI,\quad L=G_N^{-1}\Lambda^*Q_N.
 \tag{DQ1}
\]
The original frame $[I,L]$ has metric $\diag(H_K,Q_N)$: indeed
$\Lambda L=1$, $L^*G_NL=Q_N$, $I^*G_NL=0$, and the direct sum
has dimension $q$. With $J_K=H_K^{-1}I^*G_N$ write all four blocks
\[
 [I,L]^{-1}A[I,L]=\begin{pmatrix}T&C\\B&D\end{pmatrix},\quad
 B=\Lambda AI,\quad C=J_KAL,\quad T=J_KAI,\quad D=\Lambda AL.
 \tag{DQ2}
\]
In the actual support notation of DO1--9, use
$d=\max\{|\mathcal S|-1,h-\max_i h_i,1\}\le80$ (the corner term
is zero for $h=0$). The fully specified recovery map of DO6--8 satisfies
\[
 \mathcal R_{d,N}(\Lambda x,\Lambda Ax,\ldots,\Lambda A^dx)=x,
 \qquad \|\mathcal R_{d,N}\|_{Q_N^{\oplus(d+1)}\to G_N}
       \le C_{d,N}.\tag{DQ3}
\]
Here $C_{d,N}$ is exactly DO10--11, including the full corner Gram
inverses, all root differences and the original value-to-polynomial map.
It is a finite proved bound, not an additional hypothesis. One may instead
use the exact norm of the displayed recovery matrix, with the same metrics.

The feedback matrix $W_d$ is defined by the literal recursion
\[
 y_{n+1}=BT^nx+Dy_n+\sum_{a=0}^{n-1}BT^{n-1-a}Cy_a,
 \qquad y_0=0.\tag{DQ4}
\]
It sends $(BT^jx)_{0\le j<d}$ to $(y_1,\ldots,y_d)$, is lower
triangular with identity diagonal, and is invertible by subtraction.
The complete cross blocks occur in DQ4. Applying DQ3 to $Ix$ gives
\[
 \sum_{j=0}^{d-1}(T^j)^*B^*Q_NBT^j
 \succeq a_N^2H_K,
 \qquad a_N=(C_{d,N}\|W_d\|_{Q_N^{\oplus d}})^{-1}>0.
 \tag{DQ5}
\]
This follows by squaring
$\|Ix\|_{G_N}\le C_{d,N}\|W_d(BT^jx)_{j<d}\|_{Q_N^{\oplus d}}$;
thus it retains the original minimum throughout.

\section{A lower bound for multiple singular values of the actual map}
Use the metric coordinate isometries explicitly:
\[
 \widetilde T=H_K^{1/2}TH_K^{-1/2},\quad
 \widetilde B=Q_N^{1/2}BH_K^{-1/2},\quad
 \widetilde C=H_K^{1/2}CQ_N^{-1/2}.
 \tag{DQ6}
\]
These define singular values in the original metrics; no original matrix
is replaced. Put $\rho_N=\|\widetilde T\|$ and
\[
 p=\left\lceil\frac rd\right\rceil,\qquad
 S_{d,N}=\sum_{j=0}^{d-1}\rho_N^{2j},\qquad
 b_N=\frac{1}{C_{d,N}\|W_d\|\sqrt{S_{d,N}}}>0.
 \tag{DQ7}
\]
The convention $\rho_N^0=1$ includes $\rho_N=0$. In decreasing
order of singular values, the original leakage satisfies
\[
 \boxed{\sigma_p(\widetilde B)\ge b_N,
 \qquad p\ge\left\lceil\frac{8k-16}{80}\right\rceil.}
 \tag{DQ8}
\]
Here and below zero singular values are included to fill a rectangular
matrix's smaller dimension. The argument also proves that this dimension
is at least $p$.

To prove DQ8, form the exact stack
$\mathcal B_d=(\widetilde B,\widetilde B\widetilde T,\ldots,
\widetilde B\widetilde T^{d-1})^{\mathsf T}$. Equation DQ5 says
$\|\mathcal B_dv\|\ge a_N\|v\|$ for every vector $v$.
Diagonalizing $\widetilde B^*\widetilde B$ gives an orthonormal
singular system. Keep the first $p-1$ singular terms, obtaining $B_0$
with $\rank B_0\le p-1$ and
$\|\widetilde B-B_0\|=\sigma_p(\widetilde B)$.
If fewer terms exist, take $B_0=\widetilde B$ and the remaining
singular value zero. The stack with $B_0$ replacing $\widetilde B$
has rank at most $d(p-1)<r$. Choose a unit vector $v$ in its kernel.
Then
\[
 a_N^2\le\|\mathcal B_dv\|^2
 =\sum_{j<d}\|(\widetilde B-B_0)\widetilde T^jv\|^2
 \le\sigma_p(\widetilde B)^2 S_{d,N}.
 \tag{DQ9}
\]
This proves DQ8 and excludes the zero alternative. The proof uses the
whole stack and every power of the original kernel block.

If a bound without computing the feedback norm is desired, DO15 proves
$\|W_d\|\le d(2M_N)^{d-1}$, $M_N=\max(1,\|A\|_{G_N})$.
The original orthogonal compression gives $\rho_N\le M_N$, so
\[
 b_N\ge
 \bigl[C_{d,N}d(2M_N)^{d-1}
         (\sum_{j<d}M_N^{2j})^{1/2}\bigr]^{-1}.
 \tag{DQ10}
\]
Every quantity is either an actual finite norm or a previously proved
explicit bound. DQ10 makes no assertion that this scale is large enough
for a signed arithmetic asymptotic.

\section{The exact return direction and commutator}
The retained full-source Hermitian weight gives the precise identity
\[
 \widetilde C+\widetilde B^*=R_N,\qquad\rank R_N\le2,
 \tag{DQ11}
\]
proved in DO17/DC33. Take the $p$-dimensional right singular space of
$\widetilde B^*$ associated to the singular values in DQ8.
On this space $\|\widetilde B^*v\|\ge b_N\|v\|$.
Its intersection with $\ker R_N$ has dimension at least $p-2$;
there DQ11 gives $\widetilde Cv=-\widetilde B^*v$ exactly.
The min--max characterization follows directly from the eigenbasis of
$\widetilde C^*\widetilde C$: a subspace of dimension $p-2$
with Rayleigh quotient at least $b_N^2$ forces its $(p-2)$nd largest
eigenvalue to be at least $b_N^2$. Hence
\[
 \boxed{\sigma_{p-2}(\widetilde C)\ge b_N,
 \qquad \sigma_{2p-2}([A,P_{K,N}])\ge b_N,}
 \quad P_{K,N}=IJ_K.\tag{DQ12}
\]
The domain here gives $p\ge8$, so the indices are positive.
The second singular value is taken with the original $G_N$ metric
on both domain and codomain. Its isometric block presentation is
$\left(\begin{smallmatrix}0&-\widetilde C\\\widetilde B&0\end{smallmatrix}\right)$,
whose squared singular matrix is
$\diag(\widetilde B^*\widetilde B,\widetilde C^*\widetilde C)$.
Thus its singular values are the union of those of the two retained
directions, proving DQ12. The stronger multiplicity-sensitive rank
bound DC30--35 remains valid alongside these quantitative bounds.

The receiving maps are the literal $\Lambda AI$ and $J_KAL$, so
DQ8 and DQ12 give original mixed-action control. They do not set
the projection of a specified eigenclass equal to one of these singular
vectors: the actual class must still be evaluated through the same maps.
In particular the rank-two Hermitian weight in DQ11 retains cancellation
between many nonzero mixed-action directions, with their finite metric
size now bounded by DQ7--12.
\end{document}
